\section{Hàm chỉ báo là viên gạch đầu tiên}
\label{sec:ch06-ham-chi-bao}

\index{hàm chỉ báo}%
Các tiên đề đã cho thấy ràng buộc một công thức từ bên ngoài có thể quá chặt.
Ta chuyển sang trường hợp đơn giản nhất: mô hình chỉ trả lời đầu vào có nằm
trong một vùng $R$ của không gian đầu vào hay không. Vùng ấy là một hình hộp,
tức tích của $d$ khoảng, mỗi khoảng trên một tọa độ,
$R=(a_1,b_1]\times\cdots\times(a_d,b_d]$. \tn{Hàm chỉ báo}{indicator function}
của $R$ được viết

\begin{equation}
\label{eq:ch06-ham-chi-bao}
\mathbf{1}_R(y)=
\begin{cases}
1, & y\in R,\\
0, & y\notin R.
\end{cases}
\end{equation}

Với mô hình trong phương trình~\ref{eq:ch06-ham-chi-bao}, câu hỏi attribution
cụ thể hơn nhiều. Ở đầu vào $x$, đặc trưng $j$ quan trọng cho việc $x$ rơi vào
$R$ theo nghĩa nào? Bài báo đưa một câu trả lời rồi sửa nó, trên trường hợp
hai chiều~\autocite{p17firstprinciples}. Câu trả lời thứ nhất chỉ kiểm tra tọa
độ $j$ của $x$ có nằm trong khoảng $(a_j,b_j]$ của $R$ hay không: nếu có thì
$\varphi_j=1$, nếu không thì $\varphi_j=0$. Attribution nhị phân như thế không
có độ lớn, nên bài báo sửa thành câu trả lời thứ hai: giữ phép kiểm ấy nhưng
nhân thêm độ dài của khoảng còn lại. Với $d=2$, đặc trưng thứ nhất nhận
$b_2-a_2$ khi $x_1\in(a_1,b_1]$, còn đặc trưng thứ hai nhận $b_1-a_1$ khi
$x_2\in(a_2,b_2]$; với $d$ bất kỳ, sách mở rộng theo đúng cách ấy, nhân với độ
dài của mọi khoảng còn lại. Lý do là khi khoảng $(a_2,b_2]$ rộng, điều kiện
trên tọa độ thứ hai gần như luôn thỏa, nên phần lớn lý do khiến $x$ rơi vào
$R$ nằm ở tọa độ thứ nhất, và tọa độ ấy đáng nhận phần lớn attribution. Bài
báo không coi hai lựa chọn đó là cùng một attribution; chúng là hai
\tn{quy tắc nguyên tử}{atomic attribution rule} khác nhau. Chữ nguyên tử chỉ
đúng một việc: đây là attribution cho hàm chỉ báo, đơn vị nhỏ nhất mà mọi
attribution khác trong chương sẽ được ghép từ đó.

Thứ tự lập luận vì thế đi ngược chiều quen thuộc: ta quy định attribution cho
hàm chỉ báo trước, rồi mới hỏi quy định ấy mở rộng được tới lớp hàm nào.
Mục~\ref{sec:ch06-tu-nguyen-tu-den-tich-phan} làm bước mở rộng ấy, và ở đó quy
tắc nguyên tử sẽ hiện ra dưới một tên khác, là \tn{độ đo}{measure}.
